% from ega1-3-fr.tex
\oldpage[I]{107}
\begin{lemma}[3.2.6.3]
\label{I.3.2.6.3-fr}
Soient $(U_\alpha)$ un recouvrement ouvert de $X$, $(V_\lambda)$ un
recouvrement ouvert de $Y$, et supposons que pour tout couple
$(\alpha,\lambda)$, il existe un produit de $U_\alpha$ et de $V_\lambda$~;
alors il existe un produit de $X$ et $Y$.
\end{lemma}

Appliquant le lemme (\hyperref[I.3.2.6.1-fr]{3.2.6.1}) aux ouverts
$U_\alpha\cap U_\beta$ et $V_\lambda\cap V_\mu$, on voit qu'il existe un
produit des $S$-préschémas induits respectivement par $X$ et $Y$ sur ces
ouverts~; en outre, l'unicité du produit montre que, si on pose
$i=(\alpha,\lambda)$, $j=(\beta,\mu)$, il y a un isomorphisme canonique
$h_{ij}$ (resp. $h_{ji}$) de ce produit sur un $S$-préschéma $W_{ij}$
(resp. $W_{ji}$) induit par $U_\alpha\times_S V_\lambda$
(resp. $U_\beta\times_S V_\mu$) sur un ouvert~;
$f_{ij}=h_{ij}\circ h_{ji}^{-1}$ est donc un isomorphisme de $W_{ji}$ sur
$W_{ij}$. En outre, pour un troisième couple $k=(\gamma,\nu)$, on a
$f_{ik}=f_{ij}\circ f_{jk}$ dans $W_{ki}\cap W_{kj}$, comme il résulte de
l'application de (\hyperref[I.3.2.6.1-fr]{3.2.6.1}) aux ouverts
$U_\alpha\cap U_\beta\cap U_\gamma$ et
$V_\lambda\cap V_\mu\cap V_\nu$, dans $U_\beta$ et $V_\mu$
respectivement. Il y a par suite un préschéma $Z$, un recouvrement ouvert
$(Z_i)$ de l'espace sous-jacent à $Z$ et pour chaque $i$ un isomorphisme $g_i$
du préschéma induit $Z_i$ sur le préschéma
$U_\alpha\times_S V_\lambda$, de sorte que pour tout couple $(i,j)$, on ait
$f_{ij}=g_i\circ g_j^{-1}$
(\hyperref[I.2.3.1-fr]{2.3.1})~; de plus, on a
$g_i(Z_i\cap Z_j)=W_{ij}$. Si $p_i$, $q_i$, $\theta_i$ sont les projections
et le morphisme structural du $S$-préschéma
$U_\alpha\times_S V_\lambda$, on constate aussitôt que
$p_i\circ g_i=p_j\circ g_j$ dans $Z_i\cap Z_j$, et de même pour les deux
autres morphismes. On peut donc définir des morphismes de préschémas
$p:Z\to X$ (resp. $q:Z\to Y$, $\theta:Z\to S$) par la condition que $p$
(resp. $q$, $\theta$) coïncide avec $p_i\circ g_i$
(resp. $q_i\circ g_i$, $\theta_i\circ g_i$) dans chacun des $Z_i$~; $Z$,
muni de $\theta$, est alors un $S$-préschéma. Montrons maintenant que
$Z'_i=p^{-1}(U_\alpha)\cap q^{-1}(V_\lambda)$ est égal à $Z_i$. Pour tout
indice $j=(\beta,\mu)$, on a
$Z_j\cap Z'_i=g_j^{-1}(p_j^{-1}(U_\alpha)\cap q_j^{-1}(V_\lambda))$. Or,
\[
  p_j^{-1}(U_\alpha)\cap q_j^{-1}(V_\lambda)
  =p_j^{-1}(U_\alpha\cap U_\beta)\cap
   q_j^{-1}(V_\lambda\cap V_\mu)~;
\]
en vertu de (\hyperref[I.3.2.6.1-fr]{3.2.6.1}), les restrictions de $p_j$ et
$q_j$ à
$p_j^{-1}(U_\alpha)\cap q_j^{-1}(V_\lambda)$ définissent sur ce
$S$-préschéma une structure de produit de $U_\alpha\cap U_\beta$ et
$V_\lambda\cap V_\mu$~; mais l'unicité du produit entraîne alors que
$p_j^{-1}(U_\alpha)\cap q_j^{-1}(V_\lambda)=W_{ji}$. On a par suite
$Z_j\cap Z'_i=Z_j\cap Z_i$ pour tout $j$, d'où $Z'_i=Z_i$. On déduit alors
de (\hyperref[I.3.2.6.2-fr]{3.2.6.2}) que $(Z,p,q)$ est un produit de $X$ et
$Y$.

\begin{lemma}[3.2.6.4]
\label{I.3.2.6.4-fr}
Soient $\varphi:X\to S$, $\psi:Y\to S$ les morphismes structuraux de $X$
et $Y$, $(S_i)$ un recouvrement ouvert de $S$, et posons
$X_i=\varphi^{-1}(S_i)$, $Y_i=\psi^{-1}(S_i)$. Si chacun des produits
$X_i\times_S Y_i$ existe, alors $X\times_S Y$ existe.
\end{lemma}

D'après (\hyperref[I.3.2.6.3-fr]{3.2.6.3}), tout revient à prouver que les
produits $X_i\times_S Y_j$ existent quels que soient $i$ et $j$. Posons
$X_{ij}=X_i\cap X_j=\varphi^{-1}(S_i\cap S_j)$,
$Y_{ij}=Y_i\cap Y_j=\psi^{-1}(S_i\cap S_j)$~; en vertu de
(\hyperref[I.3.2.6.1-fr]{3.2.6.1}), le produit
$Z_{ij}=X_{ij}\times_S Y_{ij}$ existe. Notons maintenant que si $T$ est un
$S$-préschéma et si $g:T\to X_i$, $h:T\to Y_j$ sont des $S$-morphismes, on
a nécessairement
$\varphi(g(T))=\psi(h(T))\subset S_i\cap S_j$, d'après la définition d'un
$S$-morphisme, donc $g(T)\subset X_{ij}$ et $h(T)\subset Y_{ij}$~; il est
alors immédiat que $Z_{ij}$ est un produit de $X_i$ et $Y_j$.

\begin{env}[3.2.6.5]
\label{I.3.2.6.5-fr}
Nous pouvons maintenant achever de démontrer le th.
(\hyperref[I.3.2.6-fr]{3.2.6}). Si $S$ est un \emph{schéma affine}, il y a
des recouvrements $(U_\alpha)$, $(V_\lambda)$ de $X$ et $Y$ respectivement,
formés d'ouverts affines~; comme $U_\alpha\times_S V_\lambda$ existe en
vertu de (\hyperref[I.3.2.2-fr]{3.2.2}), il en est de même de
$X\times_S Y$ par (\hyperref[I.3.2.6.3-fr]{3.2.6.3}). Si $S$ est un
préschéma quelconque, il y a un recouvrement $(S_i)$ de $S$ formé d'ouverts
affines. Si $\varphi:X\to S$, $\psi:Y\to S$ sont les morphismes
structuraux, et si on pose $X_i=\varphi^{-1}(S_i)$,
$Y_i=\psi^{-1}(S_i)$, les produits $X_i\times_{S_i}Y_i$ existent d'après ce
qui
